=====================
====== 建置環境 ======
1. 確認node版本 : node -v
2. 在chrome安裝擴充程式React Developer Tools並釘選(如果網站是用React寫的就會亮起來)
3. 在老師的github - 用fork或下載檔案, 弄出自己的儲存庫
4. 專案資料夾絕對不能放在雲端硬碟(如iCloud的磁碟)(雲端硬碟絕對會壞掉)
5. cd至主資料夾
6. pnpm install安裝套件
7. pnpm run dev 打開Next.js
====== 建置環境 ======
=====================

-- Ai api key --
如何用api key?(sk-23d60fdcc9954737b13d2face59eb11d 老師課堂結束就會刪除)

AI api key
1. Kilo Code Ai - VS Code模組
2. 在Ai介面輸入 /settings 打開設定
3. 按供應商 - 查看更多供應商 - 輸入 DeepSeek 並點擊 - 輸入api金鑰
4. 選模型 - 預設模型 - 搜尋模型輸入 deep 選 DeepSeek 底下的 DeepSeek V4 Flash
5. 選下面的code - 一樣選DeepSeek V4 Flash (修改預設模型與Code)
6. 按右下角儲存
如果Ai對話窗出現在左邊,想移動到右邊的話:
問Ai好了
-最上面工具列
-檢視
-開啟命令選擇區 
-輸入 move 後,選 檢視:移動檢視 
-找到kilo code 
-選新增提要欄位項目

-- Ai api key --

專案裡面最重要的是app資料夾
public裡面放圖片(svg等等)

打開伺服器後在終端出現的localhost,連上瀏覽器之後選components(React的), 裡面出現的東西是 /app/page.tsx 產出的東西d

至於VsCode模組的ESlint : 這是一個會報錯抓bug的模組, 非常煩人(要習慣)
還有error lens : 容易跳出一堆錯誤的模組(可用可不用)


-- 營運部署 (一般情況下用不到這個指令)--
npm run build
打包時會跳出ESlint錯誤

npm run start
